\fsection{Ej 8 Pág 132}
\fsubsection
[\textcolor{waveBlue2}{Actividades.}]
{\textcolor{waveBlue2}{\boldit{Actividades.}}}
{
	Indica qué debemos hacer con los siguientes residuos domésticos:\\
}

\begin{solution}

	La gestión correcta de los residuos seria la siguiente:\\

	\begin{tabularx}{\linewidth}{|>{\hsize=0.2\hsize}Y|>{\hsize=1.1\hsize}Y|>{\hsize=1.7\hsize}Y|}
		\hline
		\textbf{Letra} & \textbf{Residuo}                  & \textbf{Contenedor / Gestión}                \\
		\hline
		a)             & El envase de un yogur             & Contenedor Amarillo (Envases de plástico)    \\
		\hline
		b)             & Un juguete de plástico            & Punto Limpio (No es un envase)               \\
		\hline
		c)             & Un brik de leche                  & Contenedor Amarillo (Envases mixtos)         \\
		\hline
		d)             & Un tarro de garbanzos en conserva & Contenedor Verde (Vidrio) - Tapa al Amarillo \\
		\hline
		e)             & Restos de pescado                 & Contenedor Marrón (Materia orgánica)         \\
		\hline
		f)             & El frasco de un perfume           & Contenedor Verde (Vidrio)                    \\
		\hline
		g)             & Una lavadora que ya no funciona   & Punto Limpio (Residuos RAEE)                 \\
		\hline
		h)             & Pilas                             & Contenedor específico o Punto Limpio         \\
		\hline
		i)             & Medicamentos caducados            & Punto SIGRE (Farmacia)                       \\
		\hline
		j)             & Una caja de cartón                & Contenedor Azul (Papel y cartón)             \\
		\hline
		k)             & Ropa que ya no utilizamos         & Contenedor textil o Punto Limpio             \\
		\hline
		l)             & Tapones de plástico               & Contenedor Amarillo (Envases)                \\
		\hline
	\end{tabularx}

\end{solution}



